\chapter{Reativação de registro profissional}
\label{chap:pf-reativacao}

\section{O que é?}

É a reativação de registro profissional cancelado.

\section{Quem pode utilizar esse serviço?}

Profissionais que tenham cancelado seu registro profissional e que o desejem reativar.

\section{Como solicitar o serviço?}

Por meio do \emph{site} do CRP~SP, na guia \href{https://www.crpsp.org/pagina/view/314}{Atendimento/Pessoa física}.

\section{Que informações ou documentos são necessários?}

\begin{itemize}
  \item  Documento de identificação (RG, CNH etc.);
  \item  Cadastro de Pessoa Física (CPF);
  \item  certidão de quitação eleitoral;
  \item  comprovante ou declaração de endereço;
  \item  diploma de formação em Psicologia (frente e verso);
  \item  certidão de casamento ou de união estável;
  \item  certidão de quitação militar ou certificado de reservista (para homens  até 45 anos).
\end{itemize}

\section{Quais são as etapas para a realização desse serviço?}

\begin{enumerate}
  \item  A/o psicóloga/o faz a  \href{https://www.crpsp.org/pagina/view/314}{solicitação de inscrição  prévia \emph{on-line}};
  \item  o CRP~SP faz a análise da documentação;
  \item  são emitidos os boletos para pagamento da taxa de inscrição e da  anuidade;
  \item  após o pagamento das taxas, o pedido é deferido e a declaração para  atuação profissional é enviada.
\end{enumerate}

\section{Quanto custa?}

\begin{itemize}
  \item  Taxa administrativa: R\$~105,37 (valor para 2024);
  \item  anuidade vigente proporcional (consulte a  \href{https://www.crpsp.org/pagina/view/206}{Tabela de Anuidades e  Taxas --- Pessoa Física})
\end{itemize}

\section{Quanto tempo leva?}

Até 30 dias.

\section{Como ter acesso ao atual \emph{status} do serviço?}

Ao se fazer a solicitação, é gerada uma senha de acesso. Com essa senha é possível fazer o acompanhamento \emph{on-line} do \emph{status} da solicitação, na página \href{https://cfpsp.brctotal.com/crp06_servicosonline/pgsRequerimento/Requerimentos.aspx}{Meus requerimentos}.

\section{Legislação relacionada}

\begin{itemize}
  \item  \href{https://atosoficiais.com.br/cfp/resolucao-administrativa-financeira-n-3-2007-institui-a-consolidacao-das-resolucoes-do-conselho-federal-de-psicologia?origin=instituicao&q=resolu\%C3\%A7\%C3\%A3o\%203/2007}{Resolução  CFP nº~03/2007} (Consolidação das Resoluções do Conselho Federal de  Psicologia);
  \item  \href{https://transparencia.cfp.org.br/wp-content/uploads/sites/7/2020/05/Manual-de-Procedimentos-Administrativos-e-Financeiro-CFP-RES-20.2018.pdf}{\href{https://atosoficiais.com.br/cfp/resolucao-administrativa-financeira-n-20-2018-altera-o-manual-de-procedimentos-administrativo-e-financeiro-do-sistema-conselhos-de-psicologia-anexo-da-resolucao-cfp-n-202018-a-resolucao-cfp-n-03-2007-a-resolucao-cfp-n-16-2019-e-da-outras-providencias}{Resolução  CFP nº~20/2018}}, que determina a revisão e ampliação do  \href{https://transparencia.cfp.org.br/wp-content/uploads/sites/7/2020/05/Manual-de-Procedimentos-Administrativos-e-Financeiro-CFP-RES-20.2018.pdf}{\emph{Manual  de procedimentos administrativos e financeiros}}
\end{itemize}

\section{Serviços correlatos}

\begin{itemize}
  \item  \hyperref[chap:pf-cip]{Entrega  da carteira de identidade profissional (CIP)};
  \item  \hyperref[chap:pf-inscricao_definitiva]{alteração  de inscrição provisória para definitiva};
  \item  \hyperref[chap:pf-secundaria]{inscrição secundária};
  \item  \hyperref[chap:pf-transferencia]{transferência de  registro profissional};
  \item  \hyperref[chap:pf-cancelamento]{cancelamento de  registro profissional};
  \item  \hyperref[chap:pf-prorroga_diploma]{prorrogação  do prazo de entrega do diploma};
  \item  \hyperref[chap:pf-isencao]{solicitação  de isenção e interrupção temporária de anuidades}.
\end{itemize}

\sumario